%%=============================================================================
%% Methodologie
%%=============================================================================

\chapter{\IfLanguageName{dutch}{Methodologie}{Methodology}}%
\label{ch:methodologie}

%% TODO: In dit hoofstuk geef je een korte toelichting over hoe je te werk bent
%% gegaan. Verdeel je onderzoek in grote fasen, en licht in elke fase toe wat
%% de doelstelling was, welke deliverables daar uit gekomen zijn, en welke
%% onderzoeksmethoden je daarbij toegepast hebt. Verantwoord waarom je
%% op deze manier te werk gegaan bent.
%% 
%% Voorbeelden van zulke fasen zijn: literatuurstudie, opstellen van een
%% requirements-analyse, opstellen long-list (bij vergelijkende studie),
%% selectie van geschikte tools (bij vergelijkende studie, "short-list"),
%% opzetten testopstelling/PoC, uitvoeren testen en verzamelen
%% van resultaten, analyse van resultaten, ...
%%
%% !!!!! LET OP !!!!!
%%
%% Het is uitdrukkelijk NIET de bedoeling dat je het grootste deel van de corpus
%% van je bachelorproef in dit hoofstuk verwerkt! Dit hoofdstuk is eerder een
%% kort overzicht van je plan van aanpak.
%%
%% Maak voor elke fase (behalve het literatuuronderzoek) een NIEUW HOOFDSTUK aan
%% en geef het een gepaste titel.

In dit hoofdstuk wordt de gehanteerde onderzoeksstrategie toegelicht. Gezien de technische complexiteit van de integratie tussen NVIDIA Isaac Sim en ROS 2, is gekozen voor een methodiek die de nadruk legt op systeemontwerp en een daaropvolgende diagnostische analyse van de \textit{sim-to-real gap}. Het onderzoek is onderverdeeld in vijf hoofdfasen, waarbij elke fase bijdraagt aan de beantwoording van de deelvragen binnen het probleem-, analyse- en oplossingsdomein.

\section{Fase 1: Literatuurstudie en Architectuurdefinitie}
De eerste fase richtte zich op het verkaren van theoretische kennis over de huidige stand van de techniek binnen magazijnautomatisering en de \textit{sim-to-real} problematiek. 
\begin{itemize}
    \item \textbf{Doelstelling:} Het definiëren van een hybride systeemarchitectuur waarbij de software-stack (ROS 2 Humble, Nav2) identiek blijft voor zowel de fysieke als de virtuele robot.
    \item \textbf{Onderzoeksmethode:} Literatuuronderzoek naar ROS-standaarden (REP 103/105) en vergelijkende studie van simulatie-architecturen.
    \item \textbf{Relevantie:} Deze fase legt de basis voor de beantwoording van \textbf{DV1} en \textbf{DV2} door de theoretische vereisten van coördinatenstelsels in kaart te brengen.
\end{itemize}

\section{Fase 2: Digitaal Ontwerp en Robotmodellering}
De tweede fase omvatte het mechanisch ontwerp in een iteratief proces van ongeveer vier maanden in de CAD-software Onshape. 
\begin{itemize}
    \item \textbf{Doelstelling:} Het creëren van een gedetailleerde digitale tweeling die de basis vormt voor zowel de fysieke constructie als de simulatie-import.
    \item \textbf{Onderzoeksmethode:} Modelgebaseerd ontwerpen en het gebruik van de \textit{onshape-to-robot} export-pijplijn voor URDF-generatie.
    \item \textbf{Relevantie:} De output van deze fase is direct gekoppeld aan \textbf{DV5}, waarbij de impact van export-parameters op de transformatieketen wordt onderzocht.
\end{itemize}

\section{Fase 3: Realisatie van de Fysieke Robot}
Gedurende een periode van twee maanden werd het concept uit het CAD-model omgezet naar een werkende fysieke robot. 
\begin{itemize}
    \item \textbf{Activiteiten:} Hardware-productie (3D-printen, plasmasnijden), assemblage van elektronica (Raspberry Pi 5, motordrivers, krachtbron, sensoren,...) en software-integratie van de ROS 2-stack.
    \item \textbf{Onderzoeksmethode:} Prototyping en functionele validatietesten in een echte magazijnomgeving.
    \item \textbf{Relevantie:} Deze fase dient als de "ground truth" voor \textbf{DV4}. Het succesvol mappen in de echte wereld bewijst de correctheid van de software-stack buiten de simulatie-omgeving.
\end{itemize}

\section{Fase 4: Systeemintegratie en Virtuele Implementatie}
In deze fase werd de digitale tweeling geïmporteerd in Isaac Sim en werd de communicatie met het ROS 2-domein opgezet via \textit{ActionGraphs}.
\begin{itemize}
    \item \textbf{Doelstelling:} Het realiseren van een bidirectionele bridge waarbij sensordata uit de simulator wordt verwerkt door de Raspberry Pi, en commando's worden teruggestuurd naar de virtuele motoren.
    \item \textbf{Onderzoeksmethode:} Systeemintegratietesten en het monitoren van de datacommunicatie via de \texttt{isaac\_ros\_bridge}.
    \item \textbf{Relevantie:} Deze fase adresseert \textbf{DV3} door de invloed van netwerk- en bridge-instellingen (QoS) op de datastroom te observeren.
\end{itemize}

\section{Fase 5: Diagnostische Analyse en Oplossingsstrategie}
Nadat het "rubberbanding-effect" werd vastgesteld tijdens SLAM-operaties in de simulator, verschoof de methodiek naar een diepgaand diagnostisch onderzoek.
\begin{itemize}
    \item \textbf{Doelstelling:} Het isoleren van de \textit{root-cause} van de integratiefouten en het formuleren van best-practices.
    \item \textbf{Onderzoeksmethode:} \textit{Root Cause Analysis} (RCA) door middel van stapsgewijze troubleshooting, visuele inspectie van de $TF$-tree in RViz, en log-analyse van timing-jitter en coördinaten-mismatches.
    \item \textbf{Relevantie:} Deze fase staat centraal in het beantwoorden van het analysedomein (\textbf{DV6}) en vormt de bron voor de resultaten in het oplossingsdomein (\textbf{DV7}, \textbf{DV8} en \textbf{DV9}).
\end{itemize}